The phrase \emph{protein interaction} names a family of claims organized around perturbation,
dependence, and constraint.  At the weakest level, one
part of a computation can influence another.  At a stronger level, a function contains a
nonadditive component after lower-order effects have been removed.  Stronger still, compatible
components may assemble into a probability law or effective Hamiltonian.  Structural and
biological interpretations require their own coarse graining and interventions.

This plurality is the reason a typed theory is needed.

\section{What the operator graph contributes}

The reference-centered operator graph provides a carrier that is richer than a scalar network and
weaker than a universal physical Hamiltonian.  Each residue position carries a centered
categorical fiber.  Each directed pair carries an operator from source contrasts to target
contrasts.  Support, context, interaction order, dynamics, normalization, estimation, and
validation remain additional coordinates rather than being hidden inside the edge matrix.

This carrier is broad enough to compare Potts couplings with substitution responses, and precise
enough to localize their agreement and disagreement.  It distinguishes a stable mean operator from its
background-dependent fluctuation.  It separates reciprocity of an effective pair energy from
reversibility of a route.  It shows where scalar graphs, signed contractions, and connection
Laplacians discard or preserve categorical structure.  It connects linear gauge projection to
finite probability geometry without treating a tangent coordinate as a complete joint law.

The resulting unification is coordinate-wise.  Models meet in some coordinates and remain
different in others.  That is exactly the level at which comparison becomes scientifically useful.

\section{What remains irreducibly plural}

Several pluralities are structural.

Graph identity begins with an edge criterion.  Route, statistical support, operator response,
structural contact, and experimental epistasis graphs arise from different constructions.  Their
overlap is a result.

Softmax supplies a local negative-log coordinate for whichever variable is normalized.  Site
conditionals, memory posteriors, routing distributions, and graph-state posteriors therefore carry
different local energies.  Additional compatibility and equilibrium structure support a global
Hamiltonian.

Low complexity has several axes.  Sparse support, low categorical rank, pairwise interaction order,
and a low-dimensional spectrum describe different reductions, each with its own scaling law and
approximation error.

A deep sequence model supplies a distribution of responses across backgrounds.  Its mean is one
useful static representation, while its dynamic residual records the behavior left outside that
summary.

\section{The open mathematical problems}

The framework leaves several problems deliberately unresolved.

\paragraph{From factorized edges to graph ensembles.}
Independent edge-state distributions provide local Gibbs coordinates but omit global topology.
A coupled theory must determine which degree, reciprocity, loop, geometry, and chemical constraints
belong in a graph Hamiltonian, and whether the resulting ensemble remains learnable at protein
scale.

\paragraph{From directed responses to integrable dynamics.}
Reciprocity characterizes additive pair compatibility, but deep models generate context-dependent
operators and higher-order residuals.  The appropriate integrability conditions for these fields,
and the meaning of their nonintegrable circulation, remain open.

\paragraph{From pair operators to higher-order bundles.}
Functional ANOVA identifies the mass beyond pairs, but an economical representation of typed
hyperedges must balance support, categorical rank, gauge, and realizable computation.  A pair
graph records the pairwise projection of higher-order behavior, while typed hyperedges retain the
remaining information.

\paragraph{From family inference to transferable mechanism.}
Agreement between a family Potts field and a pLM response mean would show transfer in common
coordinates.  Whether such agreement obeys a cross-family scaling law in MSA depth, alphabet,
degree, phylogenetic distortion, or model scale is an empirical question with the shape of a
statistical-physics problem.

\paragraph{From predictive operators to biological intervention.}
The most decisive evidence will preserve categorical type across multiple sequence backgrounds.
Large, linked double- and higher-order mutant panels remain rare.  Experimental design and operator
estimation must therefore be developed together; reducing assay data to a final scalar benchmark
would discard the structure the theory aims to test.

\section{The edge revisited}

Return to the pair $(i,j)$ from the opening chapter.  The structural contact, Potts block,
attention route, substitution response, and experimental epistasis may eventually converge on one
coherent account.  Such coherence comes from a sequence of explicit correspondences: common categorical coordinates,
controlled contexts, compatible directions, matched nulls, and interventions that retain the
predicted structure.

The operator graph treats unity among the five edges as a hypothesis of correspondence.  It makes
agreement precise enough to be discovered and disagreement precise enough to be informative.
That is the role of the framework, and the larger methodological claim of the book:
unification becomes stronger when translation preserves the differences that matter.
